\section{Conclusiones}\label{sec:conclusiones}

En esta práctica se han implementado y simulado con Qiskit los elementos básicos de la algoritmia cuántica, desde los oráculos reversibles hasta la estimación cuántica de fase.
En todos los ejercicios la simulación ha coincidido, dentro de la fluctuación estadística de $4096$ disparos, con la predicción teórica calculada a mano o con \texttt{Statevector}.
Las conclusiones se organizan según los objetivos de la sección~\nameref{sec:objetivos}, seguidas de las ideas que atraviesan toda la memoria y de las limitaciones del trabajo.

%--------------------------------------------------------------------
\subsection{Conclusiones por objetivo}

\paragraph{Oráculos reversibles (Bloque 1 y ejercicio 2.5).} Una función booleana que no es reversible por sí sola (el AND, o una función $\{0,1\}^4\to\{0,1\}^4$ no inyectiva)
se implementa como la permutación $U_f|x\rangle|y\rangle=|x\rangle|y\oplus f(x)\rangle$, que es unitaria y cumple $U_f^2=I$. Con ella se construyeron el AND con una puerta Toffoli,
el XOR con dos CNOT, un oráculo de 3 bits con tres puertas multicontroladas y uno de 4 bits de entrada y 4 de salida con 32 puertas. En todas las entradas de cada oráculo
(4, 4, 8, 16 y 32 entradas en los ejercicios 1.1, 1.3, 1.4, 1.5 y 1.6) la salida coincide con la tabla de verdad con probabilidad 1 y la entrada se mide intacta.
La construcción importa para el coste: el oráculo lineal $f(x)=x_4\oplus x_3\oplus x_1$ necesita 3 CNOT, frente a las 16 puertas multicontroladas que requeriría la construcción por minitérminos.
Con la entrada en superposición, el paralelismo cuántico no equivale a leer toda la tabla: cada disparo devuelve una sola de las cuatro filas, que son equiprobables
(test $\chi^2$ con $p=0{,}979$ para el AND y $p=0{,}723$ para el XOR), y hacen falta en media $8{,}31$ disparos (simulado) frente a $25/3\approx8{,}33$ (teórico) para verlas todas.

\paragraph{Retroceso de fase (Bloque 2).} Si el objetivo de una puerta controlada está en un autovector, su autovalor aparece como \emph{fase relativa} en el control, y el objetivo no cambia.
Con $CP(\pi/4)$ el qubit de control gira $\pi/4$ alrededor del eje $z$ de la esfera de Bloch, mientras que la misma puerta $P(\pi/4)$ sin control es una fase global y no se observa.
Con el CNOT y el objetivo en $|{-}\rangle$ el control pasa de $|{+}\rangle$ a $|{-}\rangle$ y se mide $1$ con probabilidad 1, de modo que la medida \emph{revela} el autovalor de $X$;
con $X$ sin control se mide $0$ en ambos casos. En el barrido de $CP(\theta)$ se obtiene $P(q_0=1)=\sin^2(\theta/2)$ con una desviación máxima de $0{,}0047$ respecto a la simulación,
y se comprueba que $CP(\theta)^k=CP(k\theta)$, que es lo que permite implementar $U^{2^j}$ en QPE. En el ejercicio 2.5, con la ancilla en $|{-}\rangle$ el oráculo booleano se convierte
en un oráculo de fase, $U_f|x\rangle|{-}\rangle=(-1)^{f(x)}|x\rangle|{-}\rangle$, y el único autovalor $-1$ corresponde a $|1\rangle|{-}\rangle$.

\paragraph{Algoritmos de Deutsch, Deutsch-Jozsa y Bernstein-Vazirani (Bloque 3).} Con una sola consulta al oráculo, Deutsch separa las funciones constantes (se mide $0$) de las balanceadas (se mide $1$)
en los cuatro oráculos de una variable, frente a las 2 evaluaciones que necesita un algoritmo clásico. En Deutsch-Jozsa con $n=2$, la función constante da $|00\rangle$ y la función XOR
da $|11\rangle$ (balanceada), con 1 consulta frente a 3 en el peor caso clásico. La función AND no es ni constante ni balanceada: el resultado es una superposición que incluye $|00\rangle$ y otros estados, de modo que el algoritmo
no puede decidir. Esto ilustra que Deutsch-Jozsa depende de la promesa sobre $f$. Bernstein-Vazirani recupera las claves $a=1010$ y $a=1101$ con un único disparo y una sola consulta, frente a las 4 de un algoritmo clásico.
En el ejercicio E1 la medida $01$ indica que la función es balanceada, y la tabla de verdad obtenida ($f=\overline{x_1}$) es coherente con el análisis como función afín, con $s=(0,1)$ y $b=1$.

\paragraph{Algoritmo de Grover (Bloque 4).} Para $N=8$ y $w=011$, una iteración lleva la amplitud del estado marcado de $-0{,}354$ a $5/(2\sqrt8)\approx0{,}884$ y su probabilidad de $0{,}125$ a $0{,}781$, tal como se calculó a mano.
El mínimo $k$ con al menos un $90\%$ de éxito es $k=2$: $P_2=0{,}9453$ teórica y $0{,}9475$ simulada, con $3872$ y $3881$ conteos para $|011\rangle$ ($\chi^2$ con $p=0{,}654$). Iterar más no mejora el resultado:
$P_3=0{,}330$, porque cada iteración gira el estado un ángulo $2\theta$ y con $k=3$ se sobrepasa el estado marcado. Alcanzar el $99\%$ no es posible en la primera subida ($k=1{,}673$ no es entero),
pero sí con $k=6$, que da $P_6=0{,}9998$ teórica y $0{,}9993$ simulada. El resultado es muy sensible a $k$ ($P_5=0{,}548$ y $P_7=0{,}577$). Con $k=2$ bastan 2 consultas al oráculo frente a las $\approx N/2=4$ de media clásicas,
pero con $k=6$ ya se necesitan más que las clásicas, de modo que para un $N$ tan pequeño solo hay ahorro de consultas en el primer caso.

\paragraph{Estimación cuántica de fase (Bloque 5).} Cuando la fase es diádica la estimación es exacta y determinista: la puerta $S$ con $t=2$ da $m=1$ ($\tilde\varphi=1/4$), y $\theta=3/32$ con $t=5$ da $m=3$ con error nulo.
Cuando no lo es, el error queda acotado por $2^{-(t+1)}$: para $\theta=1/3$ y $t=3$ se obtiene $m=3$, $\tilde\varphi=0{,}375$ y un error de $0{,}042<0{,}0625$, con $P(3)=0{,}6878$ teórica y $2826/4096\approx0{,}690$ medida.
Al aumentar $t$ el error de $\theta=1/3$ se reduce a la mitad por cada qubit ($0{,}02083$, $0{,}01042$ y $0{,}00521$ para $t=4,5,6$), y $t=6$ es el menor valor que garantiza un error inferior a $0{,}01$.
Para $\theta=2/5$ pasar de $t=2$ a $t=5$ reduce el error de $0{,}1$ a $0{,}00625$ (un factor $16$). En cambio, la probabilidad del pico \emph{no} tiende a 1 sino a $27/(4\pi^2)\approx0{,}684$, porque
$1/3$ queda siempre a $1/3$ de marca de la rejilla $m/2^t$: más qubits mejoran la estimación, pero no la probabilidad de obtener un único valor de $m$. Con el operador diagonal de dos fases ($t=6$)
el autovector $|1\rangle$ da $m=24$ sin error y $|0\rangle$ da $m=18$ con error $0{,}00625$, porque $\varphi_0=0{,}275$ no es diádica. Con una superposición $c_0|0\rangle+c_1|1\rangle$ aparecen los dos picos, $m=18$ y $m=24$,
con un peso proporcional a $|c_a|^2$ ($0{,}461$ observado frente a $0{,}458$ previsto para $m=18$ y $0{,}206$ frente a $0{,}202$ para $m=24$ con $|c_0|^2=0{,}8$), y medir el registro de conteo equivale a medir $U$:
el estado $|\psi\rangle$ colapsa sobre uno de sus autovectores.

%--------------------------------------------------------------------
\subsection{Ideas transversales}

\paragraph{El retroceso de fase como hilo conductor.} El mismo mecanismo aparece en los Bloques 3, 4 y 5: la ancilla en $|{-}\rangle$ convierte el oráculo booleano en un oráculo de fase $(-1)^{f(x)}$ en Deutsch, Deutsch-Jozsa y Bernstein-Vazirani,
el oráculo de Grover marca el estado con un signo $-1$ y QPE lee en el registro de conteo la fase del autovalor de $U$. Los ejercicios del Bloque 2, con un solo control y un solo autovector, son la versión mínima de todos ellos.

\paragraph{Concordancia entre teoría y simulación.} La Tabla~\ref{tab:conclusiones_resumen} reúne las comprobaciones cuantitativas principales. Todos los tests $\chi^2$ dan $p>0{,}05$ y las desviaciones
observadas son del orden del ruido estadístico típico $1/\sqrt{4096}\approx0{,}016$ o menores. Se ha usado \texttt{Statevector} cuando se necesitaban amplitudes y fases exactas, y \texttt{AerSimulator} cuando el resultado pedido era una medida.

\begin{table}[htbp]
    \centering
    {\small\setlength{\tabcolsep}{4pt}
    \begin{tabular}{c|p{0.40\textwidth}|c|c}
    \hline
    Bloque & Magnitud & Teórico & Simulado \\
    \hline
    1 & Disparos medios para ver las 4 filas (ej.~1.2) & $8{,}33$ & $8{,}31$ \\
    2 & $P(q_0=1)$ para $\theta=\pi/4$ (ej.~2.4) & $0{,}1464$ & $0{,}1511$ \\
    4 & $P_2$ de Grover, $N=8$ & $0{,}9453$ & $0{,}9475$ \\
    4 & $P_6$ de Grover, $N=8$ & $0{,}9998$ & $0{,}9993$ \\
    5 & $P(m=3)$ en QPE con $t=3$, $\theta=1/3$ & $0{,}6878$ & $0{,}690$ \\
    5 & Pico de QPE para $\theta=1/3$, $t=4,5,6$ & $0{,}685;\ 0{,}684;\ 0{,}684$ & $0{,}683;\ 0{,}686;\ 0{,}682$ \\
    5 & $P(m=18)$ y $P(m=24)$ con $|c_0|^2=0{,}8$ & $0{,}458;\ 0{,}202$ & $0{,}461;\ 0{,}206$ \\
    \hline
    \end{tabular}}
    \caption{Resumen de las comprobaciones teoría--simulación más relevantes.}\label{tab:conclusiones_resumen}
\end{table}

\paragraph{Ventaja cuántica y sus límites.} Los algoritmos de los Bloques 3 y 4 resuelven el problema con menos consultas que el caso clásico (1 frente a 2, 3 y 4 en Deutsch, Deutsch-Jozsa y Bernstein-Vazirani con $n=1,2,4$),
pero la ventaja depende de las hipótesis del problema (la promesa de Deutsch-Jozsa, la forma lineal de Bernstein-Vazirani) y del tamaño: en Grover con $N=8$ se pierde cuando se exige el $99\%$ de éxito.

\paragraph{Convenios de bits y reproducibilidad.} Los resultados se leen correctamente solo si se conoce el convenio de orden de los bits. Se ha seguido el convenio \textit{little-endian} de Qiskit en todos los bloques, salvo en el ejercicio 3.3 y en el E1,
donde se invierte la asignación de bits clásicos y está indicado en cada caso (sección~\nameref{sec:entorno_convenciones}). Con la semilla del simulador fijada, el cuaderno reproduce las figuras y las tablas de la memoria.

%--------------------------------------------------------------------
\subsection{Limitaciones y trabajo futuro}

Todas las simulaciones son ideales y sin ruido, por lo que resultados como la probabilidad exacta 1 de los oráculos o los picos limpios de QPE no son extrapolables a un dispositivo real.
Además, los tamaños son pequeños ($n\le4$ en los oráculos, $N=8$ en Grover, $t\le6$ en la mayor parte de QPE), y la comparación de consultas entre algoritmos cuánticos y clásicos es teórica, no una medida de tiempo.
Como continuación natural se proponen:
\begin{itemize}
    \item introducir un modelo de ruido en \texttt{AerSimulator} o ejecutar en hardware real, y medir cómo afectan la profundidad del circuito y el número de puertas multicontroladas a Grover y a QPE;
    \item extender Grover a varios estados marcados y a un número de soluciones desconocido, y ver cómo cambia el número óptimo de iteraciones;
    \item reducir el coste de QPE con variantes que usan menos qubits de conteo, como la estimación iterativa de fase, y relacionarla con aplicaciones como la factorización de Shor.
\end{itemize}
